\documentclass[12pt]{article}
\usepackage[utf8]{inputenc}
\usepackage[T1]{fontenc}
\usepackage{hyperref}
\title{Tooro Language Resources for K-12 Curriculum Development}
\author{Compiled by Claude Sonnet 5 (claude-sonnet-5) for Andr\'as Kornai}
\date{}
\begin{document}
\maketitle
\section{Introduction}
Tooro (Rutooro; ISO 639-3 ttj) is a Bantu language of western Uganda, closely related to Nyoro. Uganda's Ministry of Education defines and standardises Nyoro and Tooro as one language. Tooro is reported to lack lexical tone, which is unusual among Bantu languages. The search found a reference grammar and some orthography information, but no dictionary, primer or school textbook.
\section{Primary Educational Resources}
No Tooro primers, textbooks or school curricula turned up in the search results. This gap needs further work.
\section{Linguistic and Reference Materials}
\begin{itemize}
\item L.~T. Rubongoya, \emph{A Modern Runyoro-Rutooro Grammar} (1999, Rüdiger Köppe Verlag). English-language grammar of 334 pages. \url{https://www.koeppe.de/titel_a-modern-runyoro-rutooro-grammar-j-10}
\item Glottolog bibliographic record for Rubongoya 1999. \url{https://glottolog.org/resource/reference/id/69051}
\item Rubongoya, \emph{A Runyoro-Rutooro Grammar} (1975 typescript), held at SOAS. It includes an appendix on orthography rules. \url{https://library.soas.ac.uk/Record/PP_MS_43.01.25}
\item Wikipedia, Tooro language. \url{https://en.wikipedia.org/wiki/Tooro_language}
\item Wikipedia, Nyoro-Tooro language. It compares the 1947 and 1901 spelling systems. \url{https://en.wikipedia.org/wiki/Nyoro-Tooro_language}
\end{itemize}
\section{Religious and Community Resources}
None found in the search results.
\section{Digital and Audio Resources}
None found in the search results.
\section{Curriculum Development Context}
The language is stable and used in education, but the search found almost no published teaching materials. A curriculum would probably have to build on the Rubongoya grammar and on materials produced for the combined Nyoro-Tooro standard. Those materials should be sought from Ugandan sources.
\end{document}